\subsection{Hopfield Networks}
Given a network with $N$ neurons $\vx=[x_1, \dots, x_N]$ that take binary values, the temporal evolution dynamics of these neurons are determined by a scalar-value energy function:
\begin{equation*}
    E = - \frac{1}{2}\sum_{i,j}\omega_{ij}x_ix_j = -\frac{1}{2}\vx^\top\bm{W}\vx,\quad x_{i}, x_{j} \in \{+1, -1\}
\end{equation*}
where $\omega_{ij}$ represents the strength of connectivity between node $x_i$ and $x_j$, and the connectivity is assumed to be symmetric, \ie, $\omega_{ij} = \omega_{ji}$. We can further rewrite $\bm{W} =\sum_{i=1}^P \bm{\xi}_i \bm{\xi}_i^\top $ as a set of patterns to be stored. The update rule of each node to retrieve the most relevant pattern follows the Hebbian learning rule used in neuroscience: 
\begin{equation*}
    \vx_{t+1} = \text{sign}(\bm{W}\vx_t) = \text{sign}\left(\sum_{i=1}^P \bm{\xi}_i \bm{\xi}_i^\top\vx_t\right).
\end{equation*}

This update rule tends to minimize the energy function with retrieved patterns as its attractor. It is an embodiment of the idea of ``Neurons that fire together wire together.": If two neurons connect ($\omega_{ij}>0$), then they should have the same state ($+1$ for active and $-1$ for dead). The number of patterns the network can store and retrieve is $\mathcal{O}(N)$.

\subsection{Modern Continuous Hopfield Networks}
To overcome the limitation of linear storage capacity, modern Hopfield networks, also known as Dense Associative Memory \citep{krotov2016dense}, introduce nonlinearity in the energy and the update of neurons' states and make them suitable for continuous variables:
\begin{equation*}
    E =  -\frac{1}{2}\sum_{i=1}^Pf\left(\bm{\xi}_i^\top\vx\right),\ 
    \vx_{t+1} = \operatorname{tanh}\left(\sum_{i=1}^P \bm{\xi}_i f'\left(\bm{\xi}_i^\top\vx_t\right)\right),
\end{equation*}
where $\operatorname{tanh(\cdot)}$ is to ensure the neurons' states are constrained to the interval $[-1,1]$ so that the energy is bounded from below. Depending on the form of $f$, the network could have power or exponential storage capacity. If we set $f(x) = x^2$, this reduces to the traditional networks with linear capacity. 

If we further make modifications to the non-linearity in the energy function with $\operatorname{logsumexp(\cdot)}$, which is inspired by contrastive normalization, we can define the Modern Continuous Hopfield (MCH) energy function with a quadratic regularization term on $\vx$:
\begin{equation}
\label{eq:mch}
    E_{\text{MCH}} =  -\log\left(\sum_{i=1}^P \exp\left(\bm{\xi_i}^\top\vx\right)\right) + \frac{1}{2}\vx^\top\vx.
\end{equation}

By leveraging the concave-convex procedure \citep{yuille2003concave}, the update could be written as  
\begin{equation*}
    \vx_{t+1} =\Xi\operatorname{softmax}(\Xi^\top\vx_t),
\end{equation*}
where $\Xi = [\xi_1, \dots, \xi_P] \in \mathbb{R}^{N\times P}$. This formulation has proven to converge to stationary points of the energy function $E_{\text{MCH}}$, and is linked to the key-value memory similar to the attention mechanism \citep{ramsauer2021hopfield}. 
Notice that this update rule is essentially the cross-attention given a query vector $\vx$ and can only describe the independent evolution of that vector. It fails to faithfully cover the parallel interactions between contextual tokens in the self-attention adopted in the GPT or BERT style Transformers. 

The construction of the modern continuous Hopfield energy and update rule can also be carried out from a biologically plausible view by extending the network with hidden neurons and establishing a group of coupled differential equations. We refer the readers to \citet{krotov2021large,krotov2023new} for more details.